
\documentclass[11pt,a4paper]{article}

\usepackage[margin=2.2cm]{geometry}
\usepackage[T1]{fontenc}
\usepackage[utf8]{inputenc}
\usepackage{lmodern}
\usepackage{microtype}
\usepackage{amsmath,amssymb}
\usepackage{siunitx}
\DeclareSIUnit{\calorie}{cal}
\usepackage{booktabs,tabularx,array}
\usepackage{graphicx}
\usepackage{float}
\usepackage{caption}
\usepackage{subcaption}
\usepackage{xcolor}
\usepackage{enumitem}
\usepackage[numbers,sort&compress]{natbib}
\usepackage[colorlinks=true,linkcolor=blue!45!black,
            citecolor=green!35!black,urlcolor=blue!55!black]{hyperref}
\usepackage[nameinlink,noabbrev]{cleveref}
\crefname{figure}{Fig.}{Figs.}
\Crefname{figure}{Figure}{Figures}
\crefname{table}{Table}{Tables}
\Crefname{table}{Table}{Tables}

\graphicspath{{figures/}}
\setlength{\emergencystretch}{2em}
\setlist[itemize]{leftmargin=1.5em,itemsep=2pt,topsep=4pt}

\newcommand{\mainfigure}[5]{%
  \begin{figure}[H]
    \centering
    \IfFileExists{figures/#1}{%
      \includegraphics[width=0.98\linewidth]{#1}%
    }{%
      \fbox{\parbox[c][6.2cm][c]{0.94\linewidth}{%
        \centering\textbf{#3}\\[0.7em]
        \small #4\\[0.7em]
        \textit{Expected file: \texttt{\detokenize{figures/#1}}}}}%
    }
    \caption{\textbf{#3.} #5}
    \label{#2}
  \end{figure}%
}


% Dedicated wider placement for the dense six-panel structural plate.
% The image is centered across 1.13\linewidth, extending modestly into the
% page margins while the caption remains at normal manuscript text width.
\newcommand{\mainfigurewide}[5]{%
  \begin{figure}[H]
    \centering
    \IfFileExists{figures/#1}{%
      \makebox[\linewidth][c]{\includegraphics[width=1.13\linewidth]{#1}}%
    }{%
      \fbox{\parbox[c][6.2cm][c]{0.94\linewidth}{%
        \centering\textbf{#3}\\[0.7em]
        \small #4\\[0.7em]
        \textit{Expected file: \texttt{\detokenize{figures/#1}}}}}%
    }
    \caption{\textbf{#3.} #5}
    \label{#2}
  \end{figure}%
}

\newcommand{\SIfig}[1]{Fig.~S#1}
\newcommand{\SIfigs}[1]{Figs.~S#1}
\newcommand{\SItab}[1]{Table~S#1}
\newcommand{\SItabs}[1]{Tables~S#1}
\newcommand{\SIsec}[1]{SI Section~#1}
\newcommand{\SIsecs}[1]{SI Sections~#1}

\newcommand{\papertitle}{Matched MOF Families Reveal Regime-Dependent Pore-Geometry Effects on Adsorption}

\title{\papertitle}
\date{}

\begin{document}

\begin{center}
{\LARGE \textbf{\papertitle}}\\[1.2em]

{\large \textbf{Hosein Alimardani}$^{1,\dagger}$, \textbf{Shayan Abaei}$^{2,\dagger}$,
\textbf{Mehrdad Asgari}$^{3,*}$}\\[0.9em]

{\normalsize $^{1}$Faculty of Engineering, University of Tehran, Tehran, Iran}\\[0.35em]
{\normalsize $^{2}$School of Chemical Engineering, Faculty of Engineering, University of Tehran, Tehran, Iran}\\[0.35em]
{\normalsize $^{3}$Lucy Cavendish College, University of Cambridge, Cambridge CB3 0BU, UK}\\[0.8em]

{\normalsize $^{\dagger}$These authors contributed equally.}\\[0.35em]
{\normalsize $^{*}$Corresponding author:
\href{mailto:ma2000@cam.ac.uk}{\texttt{ma2000@cam.ac.uk}}}
\end{center}

\vspace{1em}

\begin{abstract}
Reticular chemistry makes both composition and pore architecture unusually tunable
in metal--organic frameworks (MOFs), but those variables remain strongly coupled in
large structural libraries. We therefore asked a controlled question: when major
recorded chemistry and broad architecture differences are fixed before adsorption
outcomes are inspected, what geometry-associated response remains? Using ARC--MOF,
we constructed outcome-blind matched families and froze one deterministic geometry
pair per eligible family, yielding 11,072 pairs comprising 22,144 unique
frameworks. At the CH$_4$ endpoint frozen before outcome inspection
(298 K, 65 bar), the geometry-high member showed a median uptake increase of
\SI{0.584}{\milli\mole\per\gram}
(95\% family-bootstrap interval
\SIrange{0.555}{0.615}{\milli\mole\per\gram}); 70.5\% of pairs were positive.
The response strengthened with geometry contrast, but it did not follow a universal
``larger pore is better'' rule. Accessible-surface and accessible-volume classes
were strongly favorable, whereas diameter-dominated classes had slightly negative
medians, and this separation persisted when the largest descriptor contrast was
required to dominate the remaining pore changes more clearly. The same frozen pairs
also changed response regime with operating condition: pressure activation
(28.0\%) was about 4.5 times as common as pressure suppression (6.2\%), while
capacity gains frequently opposed selectivity changes. A positive population
response survived crystallographic volumetric normalization, although strong
coupling to framework density showed that pore space and density are not independent
design knobs. Together, the matched comparisons support a chemistry-facing design
picture in which pore geometry is a multidimensional, state-dependent structural
regime rather than one scalar property to maximize. Because 99.2\% of the pairs are
hypothetical--hypothetical, this conclusion is strongest within the sampled
computational structure space and is not an experimental validation claim.
\end{abstract}

\noindent\textbf{Keywords:} metal--organic frameworks; adsorption; pore geometry;
matched families; methane storage; pressure dependence; framework density

\section{Introduction}

Metal--organic frameworks (MOFs) are compelling adsorbents precisely because
reticular chemistry allows inorganic nodes and organic linkers to be assembled into
crystalline networks whose pore size, shape, surface chemistry, and connectivity can
be varied over unusually broad ranges \cite{furukawa2013chemistry}. Isoreticular
synthesis provided an early demonstration that framework metrics can be expanded
while an underlying net is retained, with large consequences for methane adsorption
\cite{eddaoudi2002isoreticular}. Subsequent experiments have shown that methane
storage can also be altered by open metal sites, pore confinement, framework
packing, and targeted pore engineering
\cite{wu2009opensites,gandara2014al,liang2017poregeometry}. The same synthetic
tunability that makes MOFs attractive, however, complicates structure--property
interpretation: changing a linker, node, interpenetration state, or network
realization can alter chemical interactions and several aspects of pore geometry at
the same time \cite{ji2020porechemistry}.

Computational screening has made that coupling visible at a much larger scale.
Hypothetical-MOF studies associated methane capacity with pore-size distributions,
surface area, void fraction, and density \cite{wilmer2012screening,fernandez2013qspr};
computation-ready experimental collections recovered related trends across reported
structures \cite{chung2014core}; and later searches extended the design space to
hundreds of thousands or more candidate frameworks \cite{lee2021trillions}.
These approaches are indispensable for discovering high-performing regions of
materials space, and modern computational chemistry provides a mature toolkit for
combining geometric characterization with molecular adsorption simulations
\cite{coudert2016computational}. Yet a database-wide correlation answers a
different question from a matched structural comparison. Across a global library,
composition, topology, density, and pore metrics can all change together, so a
predictive descriptor need not identify what remains associated with adsorption when
major recorded chemical and architectural differences are first held fixed.

That distinction is especially important because the geometric quantities commonly
grouped under ``pore size'' describe different physical questions. Probe-accessible
surface area measures the internal framework surface that can be reached by a
finite-size probe, whereas accessible volume measures connected void space available
to that probe \cite{duren2007surface,willems2012zeopp}. Zeo++-type diameter
descriptors are different again: the largest included sphere characterizes the
largest cavity, the largest free sphere characterizes the size of a probe that can
traverse a connected pathway, and the largest included sphere along that pathway
links cavity size to network accessibility \cite{willems2012zeopp}. These global
geometric descriptors therefore do not collapse to one scalar ``bigger pore''
coordinate, nor do they specify local adsorption energies. Pore-shape analysis
likewise shows that geometrically distinct nanoporous classes can require different
optimization strategies \cite{lee2017poregeometry}, while chemistry--topology
studies demonstrate that even the benefit of a functional modification can depend on
the underlying network \cite{anderson2018poretopology}.

Adsorption thermodynamics adds a second reason to expect geometry to be
state-dependent. At low loading, occupancy is disproportionately influenced by the
strongest host--guest environments, including open metal sites and locally favorable
confinement \cite{wu2009opensites,getman2012simreview}. Stronger binding is not
automatically advantageous for storage, however, because gas retained at the
delivery state reduces usable capacity; simulations of metal-substituted MOFs and
heat-of-adsorption optimization both expose this trade-off
\cite{rana2014metal,lee2014doe}. As loading rises, the amount and organization of
accessible pore space, framework density, and adsorbate packing become increasingly
important, so high-pressure capacity and deliverable capacity occupy coupled rather
than monotonic regions of descriptor space
\cite{peng2013methane,gomezgualdron2014limits,ahmed2024methaneML,song2026poreengineering}.
A pore direction that is favorable at one state point therefore need not retain the
same value at another.

The distinction between gravimetric and volumetric performance makes that coupling
chemically and technologically consequential. A low-density framework can store a
large amount per unit mass while offering a less exceptional amount per unit crystal
or packed volume. Experimental methane studies have shown that crystallographic
density can overstate what a powder bed realizes \cite{korman2020lowpressure}, and
monolithic shaping can change the balance among pore preservation, bulk density, and
volumetric working capacity \cite{connolly2019monolith}. Process objectives add a
further layer: total uptake, working capacity, selectivity, regenerability, and
separation performance are not interchangeable targets
\cite{altintas2018co2database}. Consequently, a useful pore-geometry rule must state
the adsorbate, loading window, uptake basis, and process objective rather than
assigning one intrinsic sign to a structural descriptor.

A final challenge is sampling. MOF databases are heterogeneous across metals,
linkers, functional groups, topologies, pore geometries, and source collections, so
a pooled trend can partly reflect which regions of design space are most heavily
represented \cite{moosavi2020ecosystem}. ARC--MOF is particularly useful for a
controlled test because it combines a large multi-source framework registry with
charged CIFs, geometric descriptors, chemistry-cluster assignments, recorded
topology and dimensionality, and adsorption/process outcomes
\cite{burner2023arcmof}. Its pore-space descriptors follow established
high-throughput geometric analysis \cite{willems2012zeopp}, and its recorded
topology provides a reproducible coarse architectural label
\cite{zoubritzky2022crystalnets}. REPEAT-derived charges are part of the underlying
ARC--MOF simulation resource \cite{campana2009repeat}; they are not part of the
matching key used here. The original ARC--MOF study showed that geometric
information becomes increasingly informative toward higher loading
\cite{burner2023arcmof}. What remains unresolved is whether a recurring
geometry-associated response is still visible after major recorded chemistry and
architecture differences are removed \emph{before} adsorption outcomes are inspected.

Here we address that question with an outcome-blind matched-family design.
Frameworks are first grouped by scale-invariant whole-CIF elemental composition,
metal-element composition, and ARC--MOF metal-, linker-, and
functional-chemistry clusters; the resulting chemistry families are then refined by
recorded topology and dimensionality. Only after those matches are fixed are pore
contrasts evaluated, a 10\% geometry criterion frozen, and one deterministic pair
selected per eligible family. The primary adsorption endpoint is then opened using
the already-frozen pair identities and orientations. This design does not create a
mechanical deformation experiment and does not establish atom-by-atom identity.
Instead, it asks whether geometry-response patterns recur inside families matched on
the recorded chemistry representation and broad network architecture. The resulting
11,072 pairs support a regime-based picture: accessible-space and diameter changes
occupy different adsorption-response classes; those classes remain visible under
cleaner descriptor attribution, couple strongly to density, and change value with
loading condition and process objective. The central design implication is therefore
not to maximize a generic pore size, but to optimize a \emph{pore-geometry regime}
for a specified adsorbate and operating task.


\section{Results and Discussion}

\subsection{Outcome-blind matching defines a controlled geometry comparison}

The central design feature of the analysis is not the size of the database, but the
order in which information is allowed to enter the comparison. Recorded chemistry
was fixed first, recorded architecture second, pore geometry third, and adsorption
outcomes only after the pair catalogue had been frozen. This outcome-blind sequence
is summarized in \cref{fig:cohort}A and \SItab{1}, with the complete construction,
threshold-support, and frozen-pair accounting in \SIsecs{1--4}. ARC--MOF contained 279,019
registry records, of which 242,363 had corresponding charged CIFs that parsed
successfully. Multi-member candidate families were then defined using
scale-invariant whole-CIF elemental composition, scale-invariant metal-element
composition, and ARC--MOF metal-, linker-, and functional-chemistry cluster
assignments. Geometry, topology, dimensionality, geometric-cluster labels, partial
charges, and every adsorption or process outcome were excluded from that chemistry
key.

The chemistry match is deliberately a representation-level control rather than a
claim of atom-by-atom identity. Within those chemistry families, exact agreement in
recorded topology and dimensionality was required to define broad architectural
equivalence. Missing topology or dimensionality was treated as ambiguous rather
than imputed, and no bond reconstruction or local coordination assignment was
introduced. This refinement yielded 19,068 validated same-recorded-architecture
groups containing 44,511 frameworks, of which 18,390 retained measurable variation
in at least one carried-forward geometry descriptor (\SItab{2}--\SItab{3}). The
resulting comparison therefore asks a narrower question than database-wide
screening: whether adsorption-associated geometry patterns recur after major
recorded chemical and network differences have already been removed.

A nonzero geometric difference would have admitted many nearly indistinguishable
families, so threshold support was evaluated before outcome inspection across eight
pore descriptors: ASA, AVA, AVAf, POAVA, POAVAf, $D_i$, $D_f$, and
$D_{\mathrm{if}}$. Candidate thresholds of 5, 10, 20, and 30\% retained 14,739,
11,072, 6,948, and 5,115 families, respectively
(\cref{fig:cohort}B--C; \SItab{4}--\SItab{5}). The 10\% criterion was frozen for
the primary catalogue, while 20 and 30\% were retained as stronger prespecified
contrast tiers. The 10\% choice therefore reflects an outcome-blind support
decision rather than a threshold selected to maximize an adsorption signal.

Within every 10\%-eligible family, all possible framework pairs were compared by
the same symmetric relative-difference rule across the eight pore descriptors. The
pair with the largest maximum descriptor difference was selected, followed by the
frozen mean-difference and lexical-ID tie-breakers. One pair was retained per
family, yielding exactly 11,072 matched pairs and 22,144 unique frameworks. The
median maximum pair difference was 27.2\%, whereas the median mean difference over
all eight descriptors was 12.8\%, showing that the selected comparisons are often
multidimensional rather than clean one-variable perturbations
(\cref{fig:cohort}E). AVA was the largest-contrast class for 3,933 pairs, followed
by $D_f$ (2,211), ASA (2,036), and AVAf (1,580)
(\cref{fig:cohort}D; \SItab{6}--\SItab{7}). The catalogue also spans multiple
recorded topologies (\cref{fig:cohort}F), although pcu and fsc are the most
represented. These support patterns are part of the evidential boundary of the
study: the matched design controls recorded chemistry and broad architecture, but
it does not create an artificial series in which only one pore variable changes.

\mainfigure{Figure_1.pdf}{fig:cohort}
  {Outcome-blind construction of the matched geometry cohort}
  {(A) Cohort flow from the ARC--MOF registry to the frozen pair catalogue;
   (B) maximum within-family pore span with the frozen 10\% threshold;
   (C) family and framework support across the 5, 10, 20, and 30\% candidate
   thresholds; (D) composition of the frozen catalogue by largest-contrast geometry
   descriptor; (E) maximum versus mean pairwise geometry contrast; (F) topology
   representation among the frozen pairs.}
  {Chemistry and recorded architecture were fixed before geometry thresholding,
   deterministic pair selection, or outcome access. The frozen 10\% criterion
   retained 11,072 families and produced 11,072 matched pairs comprising 22,144
   unique frameworks. Panels D--F describe the support and structural composition
   of that already-frozen catalogue; no adsorption or process outcome enters cohort
   construction.}

For readers outside computational porous-material screening, the descriptor names
in \cref{fig:cohort} and the later response figures are translated into physical
questions in \cref{tab:descriptor-translation}. ASA and the volume descriptors are
probe-based global properties of the crystallographic void space, while the
Zeo++ diameter descriptors quantify cavity and connected-passage length scales.
Geometric accessible surface area is not identical to an experimental BET surface
area, although the two can agree closely for suitable microporous MOFs; both are
different again from a local adsorption-energy landscape
\cite{duren2007surface,walton2007bet,willems2012zeopp,thommes2015iupac}.
This distinction is important because the analysis below tests whether different
\emph{modes} of pore geometry carry different adsorption responses, not whether one
generic ``pore size'' variable should be maximized.

\begin{table}[H]
\centering
\small
\caption{Physical interpretation of the pore-geometry descriptor families used in
the matched analysis. Probe-accessible geometric surface area and its distinction
from adsorption-derived BET area follow established MOF characterization work
\cite{duren2007surface,walton2007bet}; the $D_i$, $D_f$, and
$D_{\mathrm{if}}$ definitions follow the Zeo++ sphere-based pore formalism
\cite{willems2012zeopp}. ASA and the volume descriptors are probe-based global
properties of the crystallographic void space. Density is shown for structural
context but was not included in the frozen eight-descriptor geometry threshold.}
\label{tab:descriptor-translation}
\begin{tabularx}{\linewidth}{>{\raggedright\arraybackslash}p{0.17\linewidth}
                              >{\raggedright\arraybackslash}p{0.27\linewidth}
                              >{\raggedright\arraybackslash}X
                              >{\raggedright\arraybackslash}p{0.24\linewidth}}
\toprule
Descriptor family & Geometric question & Chemistry-facing interpretation &
Interpretive boundary \\
\midrule
ASA &
What internal framework surface can a finite probe reach? &
Probe-accessible interfacial area available to an adsorbate-sized probe; changing
ASA can alter the amount of accessible framework--guest interface. &
A global geometric area, not a count of adsorption sites, an experimental BET area,
or a binding-energy measure. \\
\addlinespace
AVA / AVAf &
How much connected void space is accessible to the probe? &
Accessible volume and its fractional form describe the amount of probe-connected
pore space available within the crystal. &
Not equivalent to a characteristic pore diameter or to usable process working
capacity. \\
\addlinespace
POAVA / POAVAf &
How much probe-occupiable accessible volume is represented? &
An alternative global volume description that accounts for pore space occupiable
by a finite probe, reported in absolute/fractional source forms. &
Still a global void-space metric; it does not localize where adsorption occurs. \\
\addlinespace
$D_i$, $D_f$, $D_{\mathrm{if}}$ &
How large are the cavity and connected passage scales? &
$D_i$ is the largest included-sphere diameter, $D_f$ the largest free-sphere
diameter that can traverse a connected pathway, and $D_{\mathrm{if}}$ the largest
included sphere along that free-sphere path. &
Characteristic geometric length scales; they are not diffusion coefficients and do
not identify a unique experimental bottleneck. \\
\addlinespace
Crystal density &
How much framework mass occupies a crystallographic volume? &
A coupled structural context that links gravimetric and crystallographic volumetric
performance. &
Excluded from the frozen geometry-threshold key and not treated as an independent
causal knob. \\
\bottomrule
\end{tabularx}
\end{table}


\subsection{High-pressure methane separates accessible-space and diameter regimes}

The frozen high-pressure methane endpoint reveals a population-level preference for
the geometry-high member, but the more important result is that the sign and
magnitude depend strongly on \emph{which kind of geometry changed}. At 298 K and
65 bar, all 11,072 frozen pairs had complete CH$_4$ uptake values
(\cref{fig:primary}A--B; \SItab{8}; full primary reporting in \SIsec{5}). With the pre-existing geometry-low
$\rightarrow$ geometry-high orientation, the family-level effect was
\[
\Delta q =
q_{\mathrm{geometry\mbox{-}high}}-
q_{\mathrm{geometry\mbox{-}low}},
\]
and the median was \SI{0.5836}{\milli\mole\per\gram}. The
10,000-replicate family-bootstrap 95\% interval was
\SIrange{0.5546}{0.6149}{\milli\mole\per\gram}, and 70.5\% of pairs had
$\Delta q>0$. The corresponding mean,
\SI{2.427}{\milli\mole\per\gram}, is much larger than the median, indicating a
positively skewed distribution rather than a uniform shift of every matched family.

The response also strengthened systematically across the two stronger,
prespecified geometry-contrast tiers. The $\geq20\%$ subset
($n=6,948$) had a median of
\SI{1.698}{\milli\mole\per\gram} and was 78.3\% positive, while the
$\geq30\%$ subset ($n=5,115$) reached
\SI{3.164}{\milli\mole\per\gram} and 85.2\% positive
(\cref{fig:primary}C). The descriptive contrast bins in
\cref{fig:primary}E show the same ordering without redefining the frozen
thresholds. Geometry magnitude therefore contains information beyond the binary
10\% eligibility rule: larger matched pore contrasts are associated with larger
high-pressure methane differences. Because the bins are descriptive and the pairs
remain observational matches, this trend is evidence of response structure rather
than a mechanical dose--response experiment.

The decisive separation appears when the same primary endpoint is resolved by
largest-contrast geometry class (\cref{fig:primary}D; \SItab{9}). ASA-labelled
pairs had a median $\Delta q$ of
\SI{2.996}{\milli\mole\per\gram} and were 92.9\% positive, while AVA-labelled
pairs had a median of \SI{1.220}{\milli\mole\per\gram} and were 84.2\%
positive. AVAf and the two pore-occupiable-volume classes were also positive on
median, although POAVA and POAVAf contain much smaller populations and should not
be given the same evidential weight as ASA or AVA. In contrast, the three
diameter-labelled classes all had slightly negative medians:
$-0.146$, $-0.135$, and $-0.128$
\si{\milli\mole\per\gram} for $D_f$, $D_i$, and $D_{\mathrm{if}}$,
respectively. The pooled positive median is therefore not a universal
``geometry-high'' effect. It is the aggregate of distinct pore-space regimes with
different response directions.

This split has a straightforward chemistry-facing interpretation, provided it is
kept at the level supported by the descriptors. Accessible surface area asks how
much internal framework interface can be reached by an adsorbate-sized probe,
whereas accessible volume asks how much connected void space is available to that
probe \cite{duren2007surface,walton2007bet}. The sphere-based diameter descriptors
ask a different question: how large the cavity and connected passage scales are
\cite{willems2012zeopp}. A framework can therefore acquire a wider characteristic
passage without gaining the same amount of accessible internal surface or
probe-connected void volume, and the reverse is also possible. The three
descriptor families are correlated by crystal architecture, but they are not
different units for one scalar pore-size coordinate.

That distinction is consistent with earlier methane-storage work. Global screening
studies have repeatedly found that favorable performance occupies combinations of
surface area, pore volume or void fraction, pore size, adsorption energy, and
density rather than the maximum of any one descriptor
\cite{wilmer2012screening,fernandez2013qspr,gomezgualdron2014limits,
getman2012simreview}. Pore-shape analyses likewise show that geometrically distinct
nanoporous classes can require different optimization strategies
\cite{lee2017poregeometry}, while experimental pore engineering has demonstrated
that modifying pore geometry can materially alter methane deliverable capacity
\cite{liang2017poregeometry}. The matched-family result adds a different form of
evidence: after major recorded chemistry and broad network architecture are fixed
before outcome inspection, accessible-space-labelled pairs remain strongly
favorable while diameter-labelled populations sit near or below zero.

The physical picture should not be overextended. ASA does not count adsorption
sites, AVA does not measure binding strength, and a diameter does not identify a
unique transport bottleneck. At 65 bar, a larger inventory of accessible pore
space is plausibly relevant because substantial pore filling has occurred, but the
present database descriptors do not decompose local host--guest energies or
adsorbate packing. The appropriate conclusion from \cref{fig:primary} is therefore
not that surface area or volume ``causes'' the high-pressure response. It is that
the favorable matched-pair population is concentrated in accessible-space regimes,
whereas increasing a characteristic diameter alone is not a reliable design rule.

\mainfigure{Figure_2.pdf}{fig:primary}
  {High-pressure methane response across matched geometry pairs}
  {(A) Pair-level geometry-low versus geometry-high CH$_4$ uptake with the identity
   line; (B) 10,000 family-bootstrap medians and the primary 95\% interval;
   (C) prespecified 10, 20, and 30\% geometry-contrast tiers; (D) median response by
   largest-contrast geometry class; (E) descriptive response trend across increasing
   frozen pair contrast; (F) median response in the most represented topology
   contexts.}
  {At 298 K and 65 bar, the median matched-pair methane response is
   \SI{0.584}{\milli\mole\per\gram} (95\% family-bootstrap interval
   \SIrange{0.555}{0.615}{\milli\mole\per\gram}), with 70.5\% positive pairs.
   The response strengthens with the prespecified contrast tiers, while the
   class-resolved panel separates strongly favorable accessible-space contrasts
   from the slightly negative diameter-dominated medians. The topology panel is
   descriptive context and does not identify a unique favorable architecture.}


\subsection{Cleaner attribution and uptake-basis checks establish a coupled geometry regime}

The accessible-space/diameter separation would be less informative if it vanished
as soon as correlated pore changes were examined more strictly, or if it arose
only because uptake was expressed per unit framework mass. The post-freeze H2 and
H3 analyses test these two boundaries without changing the frozen pair catalogue,
orientation, 10\% threshold, primary endpoint, or estimand; their complete
definitions and numerical summaries are given in \SIsecs{16--17}. Together they show
that the relevant physical object is not an isolated scalar descriptor, but a
coupled pore-space regime whose high-pressure response also depends strongly on
framework density.

The first boundary is descriptor covariance. The eight pairwise pore differences
are substantially correlated (\SIfig{4}B), and the multiplicity analysis shows that
many frozen pairs exceed 10\% in several descriptors simultaneously
(\SIfig{4}D). This is chemically unsurprising in a reticular solid: changing linker
length, framework expansion, interpenetration or network realization can alter
surface-to-volume ratio, void accessibility, and characteristic length scales at
the same time. The stored largest-contrast descriptor is therefore an attribution
label for the dominant observed difference, not proof that the remaining seven
geometric dimensions were held fixed.

Within that limitation, cleaner dominance filtering preserves the central
class-specific contrast. Pooling all classes, minimum top-1/top-2 ratios of 1.10,
1.25, and 1.50 retain 4,603, 2,866, and 1,715 pairs and reduce the pooled median to
\SI{0.414}{\milli\mole\per\gram},
\SI{0.187}{\milli\mole\per\gram}, and
\SI{0.038}{\milli\mole\per\gram}, respectively
(\SIfig{4}A; \SItab{26}). Read alone, that attenuation could be mistaken for loss
of the geometry-class signal. The descriptor-resolved subsets show otherwise
(\SIfig{4}C; \SItab{27}). ASA remains strongly favorable even at a top-1/top-2
ratio of at least 1.50 ($n=311$; median
\SI{1.070}{\milli\mole\per\gram}; 90.7\% positive), while AVA remains strongly
positive at 1.25 ($n=379$; median
\SI{2.273}{\milli\mole\per\gram}; 98.2\% positive). $D_f$ remains slightly
negative at 1.50 ($n=1,265$; median
\SI{-0.087}{\milli\mole\per\gram}; 43.6\% positive). Moreover, those 1,265
$D_f$-labelled pairs constitute nearly three quarters of the 1,715 pairs retained
by the strictest pooled criterion. The near-zero pooled median is therefore partly
a population-composition shift toward the slightly unfavorable $D_f$ class, not a
uniform collapse of descriptor-specific behavior.

The class-removal analysis reaches a complementary conclusion
(\SItab{17}--\SItab{18}). Removing ASA- and AVA-labelled pairs leaves 5,103
primary-tier pairs with a median of only
\SI{0.0185}{\milli\mole\per\gram}; 51.0\% are positive and the descriptive
bootstrap interval crosses zero. Yet the same non-ASA/AVA population recovers
positive medians of \SI{0.168}{\milli\mole\per\gram} and
\SI{0.718}{\milli\mole\per\gram} in the prespecified $\geq20\%$ and
$\geq30\%$ subsets. Thus geometry \emph{mode} and geometry \emph{magnitude} are
separate pieces of the response: accessible-space changes dominate the pooled
10\% signal, while sufficiently large changes outside ASA/AVA can still be
favorable.

The second boundary is uptake basis and framework density. Re-expressing the same
frozen 65-bar pair comparison on the source volumetric basis preserves a positive
population response: the median is $+3.797$ v/v, 60.9\% of pairs are positive, and
the family-bootstrap interval is 3.438--4.185
(\SIfig{5}A--B; \SItab{28}). This rules out the narrow interpretation that the
primary sign is created solely by dividing uptake by framework mass. It does not,
however, make density irrelevant. Signed density change is strongly anticorrelated
with the gravimetric response (Spearman $\rho=-0.772$) and remains negatively
associated with the volumetric response ($\rho=-0.325$)
(\SIfig{5}C; \SItab{29}--\SItab{30}). Geometry expansion and density change are
therefore part of the same structural reorganization in many matched pairs.

The descriptor split remains visible after this normalization
(\SIfig{5}D; \SItab{31}). ASA/AVA-labelled pairs have a volumetric median of
$+11.37$ v/v and are 72.9\% positive, whereas the remaining classes have a
slightly negative volumetric median of $-0.715$ v/v and are 46.8\% positive. The
high-pressure result is consequently neither a pure gravimetric artifact nor a
density-independent effect. It is concentrated in a regime where accessible pore
space and framework mass density change together.

That coupling is well aligned with the broader methane-storage literature. High
surface area and pore volume can favor gravimetric storage, while crystal and bulk
density determine how much of that advantage survives on a volumetric basis
\cite{peng2013methane,gandara2014al,korman2020lowpressure,
connolly2019monolith}. Recent pore-engineering work likewise identifies a
nonmonotonic balance among porosity, pore volume, and framework density rather than
a benefit from maximizing any one variable \cite{song2026poreengineering}. The
matched-family analysis sharpens this point because density is not introduced as a
post-hoc predictor: it emerges as a structural quantity strongly coupled to the
already-frozen geometry direction. Regressing it away by default would therefore
risk removing part of the same physical change that the matched comparison is
designed to characterize.

The volumetric check also has a strict applicability boundary. ARC--MOF v/v is a
crystallographic/database representation; it is not a packed-bed, pellet, monolith,
or tank-level storage measurement. Experimental powders can have bulk packing
densities far below their crystallographic density, and shaping can change the
trade-off among porosity, density, and volumetric working capacity
\cite{korman2020lowpressure,connolly2019monolith}. The H3 result therefore supports
a materials-level statement---the positive accessible-space regime survives a
change of uptake basis---rather than a device-performance claim.

Taken together, \SIfig{4} and \SIfig{5} change the interpretation of
\cref{fig:primary}. The largest-contrast label is noncausal because pore
descriptors co-vary, but the accessible-space/diameter separation persists when
the dominant contrast is made cleaner. The positive population sign also survives
volumetric normalization, but its magnitude is inseparable from density changes in
many families. The defensible design object is therefore a \emph{coupled geometry
regime}: the type and magnitude of accessible pore-space change, together with the
density context in which that change occurs, rather than an isolated scalar pore
descriptor.

Topology provides a final context rather than an alternative explanation.
Removing both pcu and fsc retains 2,191 pairs with a median response of
\SI{0.295}{\milli\mole\per\gram}, 64.6\% positive, and a descriptive 95\%
family-bootstrap interval of
\SIrange{0.249}{0.351}{\milli\mole\per\gram} (\SItab{16}--\SItab{17}).
The positive sign is therefore not generated solely by the two most represented
topologies. At the same time, topology constrains pore connectivity and the
geometric ways in which reticular building blocks can be arranged, so it remains
part of the structural context rather than a nuisance category to be discarded.

\subsection{Operating condition re-orders the same geometry directions}

A pore-geometry direction that is favorable at one state point does not carry an
intrinsic sign across adsorption conditions. This is visible directly in the
redesigned discrete-endpoint comparison in \cref{fig:pressure}A--B and the complete
secondary atlas in \SItab{10}; the adsorption, process, concordance, transition,
and archetype evidence is documented in \SIsecs{6--10}. For methane storage, the median response is
essentially neutral at 5.8 bar
($-0.0047$ \si{\milli\mole\per\gram}; 49.5\% positive) but becomes
\SI{0.5836}{\milli\mole\per\gram} at 65 bar (70.5\% positive). The same direction
appears in two other CH$_4$ contexts: landfill-gas methane changes from
$-0.0004$ \si{\milli\mole\per\gram} at 0.01 bar to
\SI{0.0628}{\milli\mole\per\gram} at 4.4 bar, and natural-gas-purification
methane changes from $-0.0288$ at 0.9 bar to
\SI{0.1933}{\milli\mole\per\gram} at 9 bar. The secondary N$_2$ and H$_2$
comparisons show smaller positive shifts across their available endpoints
(\SIfig{1}), reinforcing that the magnitude of the geometry-associated response is
guest- and state-point-specific rather than a generic consequence of pair
orientation.

CO$_2$ makes the state dependence even more explicit. Landfill,
natural-gas-purification, and post-combustion CO$_2$ remain negative on median at
both stored pressure endpoints, whereas pre-combustion CO$_2$ changes from
\SI{-0.1561}{\milli\mole\per\gram} at 0.4 bar (33.2\% positive) to
\SI{0.5486}{\milli\mole\per\gram} at 16 bar (66.9\% positive)
(\cref{fig:pressure}B). The connecting segments in the figure are therefore best
read literally: they link two available database conditions. They are not
isotherms, fitted trajectories, or evidence for a continuously resolved transition
pressure.

This re-ordering is physically consistent with established adsorption chemistry.
At low loading, occupancy is disproportionately influenced by the most favorable
host--guest environments. Open-metal-site MOFs provide a direct experimental
example: neutron diffraction and adsorption measurements identify strong methane
binding at coordinatively unsaturated metal sites \cite{wu2009opensites}. Stronger
binding is not universally beneficial, however. Methane retained at the delivery
state reduces usable storage, and both metal-substitution and heat-of-adsorption
studies show that increasing interaction strength can become unfavorable under a
pressure-swing objective \cite{rana2014metal,lee2014doe}. At higher loading,
available pore inventory, packing, and density become progressively more important,
which is consistent with high-pressure screening and pore-engineering studies
\cite{gomezgualdron2014limits,liang2017poregeometry,
song2026poreengineering,getman2012simreview}. The present analysis does not decompose those
energetic and packing contributions microscopically; it shows that the \emph{same
frozen structural comparisons} change population-level sign and magnitude when the
adsorbate state point changes.

The process consequences make that state dependence more consequential than uptake
alone. Methane-storage working capacity increases by a median
\SI{0.467}{\milli\mole\per\gram} (72.3\% positive), and pre-combustion working
capacity by \SI{0.660}{\milli\mole\per\gram} (69.2\% positive)
(\cref{fig:pressure}C; \SItab{11}). Landfill, natural-gas-purification, and
post-combustion working-capacity medians are instead close to zero or negative.
Even more strikingly, selectivity decreases on median in all four separation
contexts with available values: $-0.363$ for landfill gas, $-0.638$ for
natural-gas purification, $-4.456$ for post-combustion, and $-4.705$ for
pre-combustion (\cref{fig:pressure}D). The AFM--SSP comparison in \SIfig{2}
provides the same warning from a different pair of process metrics: performance
objectives need not move together.

This divergence is chemically sensible. Working capacity rewards uptake that can
be added and subsequently released over the operating window, whereas selectivity
depends on differential affinity between mixture components. Experimental methane
studies show that high total uptake and high usable capacity are not synonymous
\cite{gandara2014al,li2014dynamic}, and large-scale separation screening likewise
finds different structural optima for selectivity, working capacity,
regenerability, and composite separation metrics \cite{altintas2018co2database}.
Chemistry--topology studies further show that the value of a functional
modification can depend on the network in which it is expressed
\cite{anderson2018poretopology}. The negative selectivity medians in
\cref{fig:pressure}D should therefore not be interpreted as a contradiction of the
high-pressure methane result. They demonstrate that \emph{capacity benefit and
separation benefit are different materials questions}.

\mainfigure{Figure_3.pdf}{fig:pressure}
  {Operating condition changes the direction and process value of the matched
   geometry response}
  {(A) Discrete lower- and higher-pressure CH$_4$ endpoints for methane storage,
   landfill gas, and natural-gas purification; (B) corresponding CO$_2$ endpoints
   for the four separation contexts; (C) median working-capacity change;
   (D) median selectivity change. The complete adsorption and process summaries are
   reported in \SItab{10} and \SItab{11}, with secondary N$_2$/H$_2$ and AFM--SSP
   views in \SIfig{1} and \SIfig{2}.}
  {The same frozen pairs can be neutral or unfavorable at lower loading and
   favorable at higher loading, while the process consequence depends on the
   objective being optimized. Methane-storage and pre-combustion working capacities
   increase on median, whereas selectivity decreases in all four separation
   contexts with available values. The connecting segments in A and B join only
   the stored discrete pressure endpoints and do not define a continuous crossover
   pressure.}


\subsection{Pressure activation is asymmetric and geometry-class dependent}

Cross-condition concordance shows that the 65-bar methane response is not simply a
scaled version of every other adsorption endpoint. Several higher-loading
conditions correlate positively with the primary effect---for example,
$\rho=0.777$ for pre-combustion CO$_2$ at 16 bar and $\rho=0.640$ for
natural-gas-purification CH$_4$ at 9 bar---whereas several low-pressure CO$_2$
conditions correlate negatively ($\rho\approx-0.28$ to $-0.32$)
(\SIfig{3}; \SItab{12}). Rank correlation and sign agreement answer different
questions, but both show that operating condition reorganizes the pair-level
response rather than merely changing its scale.

The low-to-high sign transitions make that asymmetry intuitive
(\SItab{13}). For methane storage, 26.8\% of pairs move from negative at 5.8 bar
to positive at 65 bar, compared with only 5.8\% in the opposite direction.
Landfill and natural-gas-purification CH$_4$ show similar negative-to-positive
fractions of 26.6 and 26.8\%, respectively. Pre-combustion CO$_2$ is the most
pronounced example: 43.5\% of pairs switch from negative at 0.4 bar to positive at
16 bar, whereas 9.8\% switch from positive to negative. These are comparisons
between designated discrete endpoints, not kinetic switching events or
framework-transition measurements.

The pair-level summary in \cref{fig:archetypes}A converts this condition dependence
into four deterministic majority-sign regimes. The 50\% boundaries separate
positive-both, pressure-activated, pressure-suppressed, and negative-both pairs.
Among all 11,072 frozen pairs, 3,618 (32.7\%) are positive in both regimes, 3,100
(28.0\%) are pressure activated, 687 (6.2\%) are pressure suppressed, and 3,667
(33.1\%) are negative in both (\cref{fig:archetypes}B; \SItab{14}). Pressure
activation is therefore approximately 4.5 times as common as pressure suppression,
and the median fraction of positive endpoints rises from 44.4\% at lower pressure
to 55.6\% at higher pressure. This population asymmetry is one of the clearest
signatures that the geometry response is state-dependent rather than a fixed
descriptor property.

Crucially, the archetypes are not distributed randomly across geometry modes
(\cref{fig:archetypes}C; \SItab{15}). ASA-labelled pairs are strongly enriched in
favorable behavior: 57.4\% are positive in both regimes and another 34.9\% are
pressure activated. AVA-labelled pairs are especially enriched in pressure
activation (40.4\%). In contrast, more than half of the $D_f$-, $D_i$-, and
$D_{\mathrm{if}}$-labelled populations remain negative in both regimes. This is
the pressure-domain counterpart of \cref{fig:primary}D: accessible-space and
diameter contrasts are not only different at the single 65-bar endpoint; they
occupy different \emph{response-regime compositions} across the wider endpoint
inventory.

Topology changes those frequencies but does not replace the geometry-class picture.
Among topology classes represented by at least 100 frozen pairs, fsc has a higher
pressure-activated fraction than pcu, nbo, or sra
(\cref{fig:archetypes}D; \SItab{16}). Yet pressure activation occurs across
multiple topologies, and removing the two dominant pcu/fsc classes does not remove
the positive primary sign (\SItab{17}). The evidence is therefore better described
as topology-modulated geometry response than as a single topology mechanism.
Topology constrains connectivity and the geometrical ways in which nodes and
linkers fill three-dimensional space, but the data do not identify one network as
the universal origin of pressure activation.

\mainfigure{Figure_4.pdf}{fig:archetypes}
  {Pressure-response archetypes expose geometry-class heterogeneity}
  {(A) Pair-level fraction of positive lower- versus higher-pressure endpoints, with
   the 50\% majority-sign boundaries defining four response regions; (B) overall
   prevalence of the four archetypes; (C) archetype composition by largest-contrast
   geometry class; (D) pressure-activated fraction for topology classes represented
   by at least 100 frozen pairs. Numerical archetype, class, and topology summaries
   are reported in \SItab{14}--\SItab{16}.}
  {Pressure activation is substantially more common than pressure suppression
   (28.0 versus 6.2\% of pairs; approximately 4.5-fold). ASA- and AVA-labelled
   pairs are enriched in positive-both and pressure-activated regimes, whereas
   diameter-dominated pairs more often remain negative across both pressure groups.
   Topology modulates prevalence but does not define a unique favorable
   architecture or a causal activation mechanism.}


\subsection{Representative crystal structures translate descriptor regimes into
crystal-space intuition}

The statistical descriptor classes become chemically more interpretable when the
underlying framework records are viewed as matched structures rather than as rows
in a table. Six frozen pairs were therefore selected deterministically from
predeclared geometry and response strata for structural visualization
(\cref{fig:structures}; complete definitions and provenance in \SItab{33} and
\SIsec{19}).
Selection was mutually exclusive across panels, did not maximize adsorption
effect, and did not use visual appeal. The 12 CIFs were extracted with checksums and
rendered using a common convention within each pair. Figure 5 is consequently not
a gallery of unusually attractive structures; it is a crystallographic
translation of six predeclared population regimes.

\mainfigurewide{Figure5_Structural_Panel.pdf}{fig:structures}
  {Representative matched structures across geometry-response regimes}
  {(A) Broad accessible-space contrast: pcu, pressure activated and ordered by ASA,
   with multiple accessibility descriptors changing together; (B) AVA-led
   accessible-space contrast: fsc, pressure activated; (C) ASA--diameter mismatch:
   fsc, positive in both pressure regimes; (D) diameter counterexample: pcu,
   negative in both regimes and ordered by $D_f$; (E) pressure-suppressed diameter:
   pcu, with co-dominant $D_i$/$D_{\mathrm{if}}$ changes; (F)
   experimental-source anchor: dia, positive in both regimes, with co-dominant
   AVA/AVAf. Full case definitions and source metrics are given in \SItab{33}.}
  {Each panel places the frozen geometry-low and geometry-high structures side by
   side and uses a local inset only for crystallographic context. Panels A--C anchor
   favorable accessible-space regimes, whereas D and E show that increasing a
   characteristic pore diameter does not guarantee a favorable adsorption response.
   Panel F supplies an experimental-source crystallographic record, not experimental
   validation of the simulated adsorption response. The boxed regions are not
   adsorption-site or unique-bottleneck assignments, and the global descriptors can
   co-vary within every pair.}

Panels A--C show three chemically distinct ways in which a pair can enter the
favorable accessible-space population. Panel A is the most dramatic but also the
least isolated: its frozen ordering label is ASA, yet several accessibility
descriptors move from an inaccessible/zero state to substantial nonzero values.
The correct visual interpretation is therefore a \emph{broad accessible-space
opening}, not an ASA-only perturbation. Its large positive
65-bar response (\SI{9.10}{\milli\mole\per\gram}) is consistent with the
population enrichment of accessible-space classes, but the structure cannot assign
that response to one surface element or local site. The descriptor-correlation and
multiplicity analyses in \SIfig{4} are essential context for reading this panel.

Panel B occupies a related but less extreme region: AVA is the largest contrast,
the pair is pressure activated, and multiple pore descriptors increase together.
It visualizes a volume-led accessible-space change rather than a clean AVA
experiment. Panel C is more chemically revealing. ASA increases from 180.313 to
353.390 in the stored source units, while AVA is essentially unchanged
(268.519 to 267.599) and $D_f$ decreases from 8.495 to 5.961~\AA. The favorable
\SI{2.67}{\milli\mole\per\gram} response therefore accompanies \emph{more
accessible internal surface without a parallel increase in accessible volume or
characteristic free-sphere diameter}. This single structural example does not
establish mechanism, but it makes the central descriptor distinction tangible:
surface accessibility, connected void volume, and characteristic diameter encode
different geometrical aspects of a pore network.

Panels D and E supply the complementary structural counter-regimes. In D, $D_f$
increases from approximately 9.51 to 11.51~\AA\ while crystal density is essentially
unchanged, yet the 65-bar methane response is
\SI{-0.460}{\milli\mole\per\gram} and the pair remains negative in both pressure
regimes. In E, $D_i$ and $D_{\mathrm{if}}$ increase together
($D_i\approx8.89\rightarrow10.72$~\AA), density again changes negligibly, and the
pair is pressure suppressed with a primary response of
\SI{-0.188}{\milli\mole\per\gram}. These cases are particularly useful because
they separate the intuitive idea of a wider passage/cavity from the idea of a more
favorable adsorption environment. Experimental methane work provides a parallel
warning: deliberately filling pore space and reducing pore size in an IRMOF-8
family improved volumetric uptake and deliverable capacity by changing both
confinement and adsorption interactions \cite{tran2015porefilling}, while targeted
mesopore engineering can be beneficial in a different high-pressure regime
\cite{liang2017poregeometry}. The chemically meaningful lesson is therefore not
``small pores are better'' or ``large pores are better'', but that the
\emph{organization of usable pore space relative to the adsorption state point}
matters.

Panel F plays a different role. Its dia pair comes from an
experimental-source structural collection, has an AVA-led/co-dominant accessible
space contrast, and shows only a small positive primary response
(\SI{0.083}{\milli\mole\per\gram}). Its inclusion is valuable precisely because it
does not maximize effect size: it anchors the visualization in crystallographic
records associated with experimental-source structures. The adsorption values,
however, remain simulated ARC--MOF outcomes. Panel F is therefore evidence of
structural provenance, not experimental adsorption validation.

Taken as a whole, \cref{fig:structures} connects the statistical result to a
chemistry-facing picture without raising the evidential level. A pore network can
change its surface-to-volume character, accessible void inventory, and
characteristic passage dimensions in different proportions. Reticular chemistry
makes such coupled changes natural because changing building-block dimensions,
interpenetration, or network realization can reorganize several geometric
properties at once \cite{furukawa2013chemistry,ji2020porechemistry}. The Figure 5
insets help the reader see representative pore/channel environments, but they are
not maps of adsorption density, measured apertures, or local energetic minima.
Their purpose is to show why the global descriptor families in
\cref{tab:descriptor-translation} and \cref{fig:primary} should not be collapsed
into one visual notion of ``pore size''.


\subsection{Source heterogeneity defines an applicability boundary, not a mechanism}

ARC--MOF combines several source collections that sample chemical and geometric
design space unevenly. Such heterogeneity can, in principle, make a pooled trend
depend on whichever source occupies the most favorable part of that space
\cite{chung2014core,moosavi2020ecosystem}. Broad provenance and the post-freeze
same-source analysis are documented in \SIsecs{12 and 18}. The leave-source-out results address
that possibility at the level the data permit: removing each major source in turn
retains a positive primary median, and the ASA/AVA population remains positive
after removal of DB0 or DB5 (\SItab{20}--\SItab{21}). No single major source is
therefore required for the positive population sign.

Source-specific effect size nevertheless varies strongly, and \SIfig{6} shows why
the variation should be interpreted as sampled-space heterogeneity rather than a
database-source mechanism. Among same-source populations with at least 100 pairs,
DB5 has a median primary response of
\SI{10.437}{\milli\mole\per\gram} and 99.4\% positive pairs, whereas DB6 has a
median of \SI{-0.244}{\milli\mole\per\gram} and 34.2\% positive
(\SItab{32}). Those populations differ simultaneously in geometry magnitude and
class composition: the median maximum pair contrast is 0.958 for DB5 and 0.180 for
DB6, while ASA/AVA-labelled pairs constitute 72.2\% and 8.0\%, respectively.
Source identity is therefore entangled with the structural regimes being sampled.
It is useful as a provenance/applicability label, but not as an explanatory causal
variable.

The more restrictive generalization boundary is experimental coverage
(\SItab{19}). Of the 11,072 frozen pairs, 10,981 (99.2\%) are
hypothetical--hypothetical, 82 are experimental--experimental, and 9 are mixed.
The experimental--experimental subset is directionally consistent with the
primary response (median \SI{0.355}{\milli\mole\per\gram}; 78.0\% positive), but
its small size cannot support broad experimental transferability. The present
evidence is therefore strongest as a statement about recurrent geometry-response
regimes within the sampled ARC--MOF structure space. Experimental testing of
purposefully chosen matched families remains the next route to determining how
far those regimes generalize to synthesized, activated materials.


\subsection{A regime-specific design rule emerges from the matched families}

The combined evidence in
\cref{fig:primary,fig:pressure,fig:archetypes,fig:structures} rejects a
one-dimensional pore-maximization rule. A compact claim--evidence--boundary map is
provided in \SItab{22}. At the primary high-pressure methane
condition, accessible-area and accessible-volume classes are strongly enriched in
favorable responses; the distinction from diameter classes survives cleaner
descriptor dominance (\SIfig{4}) and remains visible after volumetric normalization
(\SIfig{5}). The same structural direction can then move toward or away from
favorable behavior as the adsorbate, pressure window, and process metric change.
Figure 5 makes that statistical result crystallographically intuitive: ``geometry''
is not one degree of freedom.

A useful design description therefore needs at least five coordinates. The first is
\emph{geometry mode}: accessible surface, accessible volume, and characteristic
diameter answer different structural questions. The second is \emph{geometry
magnitude}: the 20 and 30\% prespecified contrast tiers show that larger matched
changes can matter even outside the ASA/AVA population. The third is
\emph{structural coupling}: pore descriptors co-vary and the high-pressure response
is strongly linked to crystal-density change. The fourth is the \emph{state point}:
low- and high-loading endpoints do not preserve the same sign landscape. The fifth
is the \emph{objective}: total uptake, working capacity, selectivity, and other
process metrics reward different combinations of adsorption and release.

This five-coordinate view reconciles several strands of methane-storage literature.
Large-scale screening has long placed good methane performance in particular
combinations of void fraction, surface area, pore size, adsorption energy, and
density \cite{wilmer2012screening,fernandez2013qspr,
gomezgualdron2014limits}. Pore-shape analyses show that different geometric classes
can require different optimization strategies \cite{lee2017poregeometry};
experiments demonstrate that both targeted pore enlargement
\cite{liang2017poregeometry,song2026poreengineering} and deliberate pore filling
\cite{tran2015porefilling} can be beneficial in different structural and operating
contexts. Density and shaping further determine how gravimetric promise translates
toward volumetric performance
\cite{peng2013methane,korman2020lowpressure,connolly2019monolith}. The matched
families add a complementary evidential layer: these distinctions recur after major
recorded chemistry and broad architecture have been fixed before the primary
outcome is inspected.

The resulting design rule is deliberately conditional:

\begin{quote}
\emph{Optimize a pore-geometry regime for a specified adsorbate, loading window,
density context, and process objective; do not maximize one scalar pore descriptor
in isolation.}
\end{quote}

For the present matched population, accessible-space expansion is the dominant
favorable regime for high-pressure methane capacity. That statement does not imply
that every ASA/AVA increase is favorable, that diameter should be minimized, or
that one can independently tune these descriptors in a real synthesis. It means
that, under the recorded chemistry/architecture controls and within the sampled
structure space, the recurring response is organized by geometry mode and operating
context more strongly than by a universal ``larger-pore'' direction.


\subsection{Scope and limitations}

The strength of the matched-family design is also the source of its principal
boundary: it controls a recorded representation, not every microscopic chemical
degree of freedom. Whole-CIF elemental composition, metal-element composition,
ARC--MOF chemistry-cluster assignments, recorded topology, and dimensionality are
held fixed, but no bond graph, oxidation state, linker identity, node identity, or
local coordination environment is reconstructed. The matched pairs are therefore
not controlled mechanical deformations of a single crystal, and the
largest-contrast descriptor remains an attribution label within correlated pore
changes rather than a causal perturbation.

The adsorption evidence inherits the scope of ARC--MOF
\cite{burner2023arcmof}. The reported values are equilibrium simulation outcomes
for recorded structures; they do not include synthesis probability, activation
success, defects, phase mixtures, structural breathing not represented by the
record, kinetic transport, or long-term stability. Classical interaction models and
assigned charges introduce force-field and electrostatic uncertainty that is not
propagated through the family bootstrap. The database-specific limitations and full
interpretive boundary are documented in \SIsecs{20 and 22}, with the reproducibility
audit in \SIsec{21}.

Operating-condition conclusions are similarly bounded by the available data grid.
``Pressure activated'' and ``pressure suppressed'' classify changes in majority
sign between designated lower- and higher-pressure endpoint sets; they do not
identify a continuously resolved crossover pressure or a kinetic pressure-induced
transition. The process metrics are equilibrium screening quantities rather than
complete cyclic-process simulations. Likewise, the v/v robustness result is a
crystallographic/database uptake-basis check, not packed-bed, pellet, monolith, or
tank-level deliverable performance.

Finally, transferability to experiment remains the largest open question. Because
99.2\% of the frozen pairs are hypothetical--hypothetical, Figure 5F and the small
experimental--experimental subset provide provenance context rather than validation.
The most informative experimental follow-up would be a deliberately chosen set of
chemistry- and architecture-matched frameworks spanning accessible-space and
diameter regimes, measured at multiple loading conditions and reported on both
gravimetric and physically relevant volumetric bases. Such a study would test the
regime picture directly without changing the evidential status of the present
computational result.

\section{Methods}

\subsection{Data resources, identifier harmonisation, and structural parsing}

The analysis used the ARC--MOF framework registry, charged CIF archive, geometric
descriptors, topology and dimensionality assignments, RAC/chemistry-cluster tables,
adsorption records, and process tables
\cite{burner2023arcmof,willems2012zeopp,janet2017rac}. A canonical framework
registry was constructed before matching. Source identifiers were normalized by
explicit deterministic rules rather than row order, and auxiliary sources were
joined through the canonical identifier. Duplicate keys, normalization collisions,
unmatched identifiers, source coverage, and required schemas were audited before
the corresponding stage was treated as complete.

Charged CIFs were parsed with pymatgen \cite{ong2013pymatgen} without primitive-cell
reduction. The extraction recorded formula and elemental composition, normalized
element fractions, metal-element composition, site count, cell volume and density,
disorder/partial-occupancy information, oxidation-state evidence where explicitly
present, and raw charge-tag metadata. Metal classification referred only to
chemical element class and did not assert framework-node membership. Oxidation
state, framework-connected linker identity, local coordination, bonds, and nodes
were not inferred from the CIFs.

\subsection{Recorded-chemistry and architecture matching}

The recorded-chemistry signature contained five components: scale-invariant
whole-CIF elemental composition, scale-invariant metal-element composition, and the
ARC--MOF metal-, linker-, and functional-chemistry cluster assignments. Eligibility
required a successful and internally consistent CIF parse, an ordered structure,
at least one metal element, complete RAC evidence, and valid nonnegative
chemistry-cluster labels. Negative cluster assignments were treated as unique and
were never grouped. Mixed-metal structures were retained only when normalized
whole-CIF composition, metal-element composition, and all three chemistry-cluster
assignments matched. Geometry, topology, dimensionality, geometric-cluster
identity, partial charges, adsorption outcomes, and process outcomes were excluded
from the chemistry signature.

Architecture was refined only after chemistry families were fixed. The architecture
key consisted of chemistry-family identifier, recorded topology, and recorded
dimensionality. Missing topology or dimensionality was treated as ambiguous rather
than imputed, and measured geometry did not enter the architecture key. No CIF was
reparsed to infer a replacement network label. Recorded topology was therefore used
as a practical coarse architectural descriptor rather than as proof of complete
atom-by-atom structural identity \cite{zoubritzky2022crystalnets}.

\subsection{Geometry threshold and deterministic pair construction}

Ten structural quantities were retained descriptively: unit-cell volume, density,
ASA, AVA, AVAf, POAVA, POAVAf, $D_i$, $D_f$, and $D_{\mathrm{if}}$.
Scientific thresholding used only the eight pore descriptors, excluding unit-cell
volume and density. For descriptor $x$, the within-family fractional span was
\[
\Delta_{\mathrm{rel}}(x)=
\frac{x_{\max}-x_{\min}}
{\max\left(|x_{\min}|,|x_{\max}|\right)}.
\]
Candidate thresholds of 5, 10, 20, and 30\% were evaluated without reading
adsorption or process outcomes. The 10\% threshold was frozen as the primary
geometry criterion before pair construction; the 20 and 30\% thresholds were
retained as stronger-contrast sensitivity tiers.

Within each 10\%-eligible family, every unordered framework pair was enumerated.
For pore descriptor $k$, the pairwise relative difference was
\[
d_k(i,j)=
\frac{|x_{ik}-x_{jk}|}
{\max\left(|x_{ik}|,|x_{jk}|\right)},
\]
with zero assigned when both values were zero. The selected pair maximized
$D_{\max}(i,j)=\max_k d_k(i,j)$. Ties were resolved by the larger mean
eight-descriptor difference and then lexicographically by framework identifier.
Exactly one pair was frozen per eligible family. The descriptor attaining
$D_{\max}$ was retained as the pair's largest-contrast geometry label and later
used to orient the pair from its lower to its higher descriptor value. No outcome
entered pair selection or orientation.

\subsection{Outcome inventory, primary endpoint, and uncertainty}

Before adsorption values were inspected, outcome schemas and framework coverage
were inventoried using identifier and condition metadata only. The confirmatory
endpoint was then frozen as CH$_4$ uptake at 298 K and 65 bar in
\si{\milli\mole\per\gram}. A pair contributed only if both members had observed
values; no imputation was used. The primary signed effect was
\[
\Delta q=q_{\mathrm{high}}-q_{\mathrm{low}},
\]
where high and low refer to the already-frozen orientation of the pair-specific
largest-contrast geometry label. The population estimator was the median
$\Delta q$ across frozen families. Uncertainty was estimated by nonparametric
percentile bootstrap, resampling whole matched families 10,000 times with fixed
seed 17. This interval quantifies variation across matched families and does not
propagate the per-framework Monte Carlo sampling uncertainty of the source
simulations.

Supporting same-endpoint summaries included the mean signed difference,
positive-pair fraction, and signed symmetric relative effect,
\[
r_{\mathrm{sym}}=
\frac{2(q_{\mathrm{high}}-q_{\mathrm{low}})}
{|q_{\mathrm{high}}|+|q_{\mathrm{low}}|},
\]
with zero assigned when both values were zero. The 20 and 30\% geometry tiers were
prespecified sensitivity populations. No minimum practically meaningful uptake
difference was prespecified.

\subsection{Secondary adsorption, process, and cross-condition analyses}

After completion and audit of the primary analysis, the remaining adsorption and
process outcomes were treated as secondary or exploratory while preserving pair
membership and orientation. Adsorption endpoints included
\si{\milli\mole\per\gram}, v/v, wt\%, heat of adsorption, and mixture selectivity
where available. Process endpoints included uptake, working capacity, selectivity,
purity, SSP, and AFM. For every secondary endpoint the signed effect used the same
frozen direction, $\Delta=y_{\mathrm{high}}-y_{\mathrm{low}}$. These analyses
were descriptive rather than a family of confirmatory hypothesis tests; no
secondary endpoint was promoted to confirmatory status after value inspection and
no multiplicity-adjusted significance claim was made.

For cross-condition analysis, the 18 \si{\milli\mole\per\gram} adsorption
endpoints were assembled into one pair-level effect matrix. Each endpoint was
compared with the primary 65-bar CH$_4$ response by Spearman correlation, sign
agreement, and positive fraction. Within each source/guest series, the lower- and
higher-pressure endpoints were also compared directly. Across the nine lower- and
nine higher-pressure endpoints, a pair was defined as majority positive within a
regime when strictly more than half of its observed effects were positive. Four
deterministic response archetypes were then assigned: positive in both regimes,
pressure activated, pressure suppressed, and negative in both. Geometry-class- and
topology-specific archetype frequencies were computed without redefining pair
membership.

\subsection{Robustness, provenance, and structural case selection}

Topology- and largest-contrast-class removal analyses preserved the frozen pair
catalogue, orientation, endpoint, and estimand. Geometry-attribution stability
ranked the eight pairwise pore-descriptor differences within every frozen pair and
compared the largest with the second largest. Top-1/top-2 dominance cuts of 1.10,
1.25, and 1.50 were applied without reselecting pairs, and summaries were examined
both pooled and by largest-contrast geometry class.

Uptake-basis robustness first required exact reproduction of the frozen
\si{\milli\mole\per\gram} pair effects and then evaluated the same orientation
using the source v/v and wt\% fields. Density-effect correlations, density-quintile
summaries, and ASA/AVA versus non-ASA/AVA volumetric summaries were treated as
descriptive robustness checks, not as new matching variables or confirmatory
estimands. Density was therefore not retrospectively added to the matching key or
regressed away after outcome inspection.

ARC--MOF source provenance was recovered from the \texttt{DBx-} prefix encoded in
framework identifiers following the source mapping of the database
\cite{burner2023arcmof}. Each frozen pair was classified as same-source or
cross-source and by hypothetical/experimental source category. For source
populations represented in at least 100 frozen pairs, leave-one-source-out and
same-source summaries were used as descriptive applicability tests. Source labels
were not interpreted causally because database provenance is correlated with
structure-generation protocol, chemical space, and geometry distributions.

Six structural examples were selected deterministically from predeclared
response/geometry strata only after the population analyses were frozen. No pair
was reused across panels. The corresponding 12 CIFs were extracted with SHA256
checksums for visualization; these examples do not alter or supplement the
population-level estimands.

\subsection{Audit and reproducibility}

Every completed stage produced deterministic framework-, family-, pair-, or
endpoint-level tables together with provenance manifests. Independent audit scripts
reconstructed grouping keys, exclusions, threshold flags, pair optimality, outcome
linkage, bootstrap summaries, cross-condition matrices, archetype assignments,
robustness subsets, source labels, and selected structural cases. The audit verified
the temporal ordering central to the design: chemistry and architecture were fixed
before geometry thresholding; pair membership and orientation were fixed before
outcome access; the primary endpoint was frozen before its values were inspected;
and all later analyses preserved those choices. Full procedural detail, closure
tables, and secondary results are reported in the Supporting Information. The
claim--evidence map is \SItab{22}; the complete secondary adsorption and process
closures are \SItab{23}--\SItab{24}; and the source/code map is \SItab{25}
(\SIsecs{13--15 and 21}).

\section{Conclusions}

By constructing 11,072 outcome-blind matched MOF pairs, this study asks a controlled
question that conventional database-wide screening does not: what pore-geometry
response remains after major recorded chemistry and broad architectural differences
are fixed before the primary adsorption values are inspected? At 298 K and 65 bar,
the answer is a clear positive population response for methane, with a median
$\Delta q$ of \SI{0.584}{\milli\mole\per\gram} and 70.5\% positive pairs, but the
average conceals strong structural heterogeneity.

The decisive result is not simply that ``geometry matters''. Accessible-surface-area
and accessible-volume regimes are strongly favorable, whereas diameter-dominated
regimes have slightly negative medians at the primary 10\% tier. That separation
survives cleaner top-1/top-2 attribution filters, showing that it is not erased when
the largest geometry contrast is more clearly distinguished from the other pore
descriptors. At the same time, the pooled response outside ASA/AVA approaches zero
at 10\% but becomes increasingly favorable at the 20 and 30\% contrast tiers, so
geometry mode and geometry magnitude must be considered separately.

The uptake-basis and pressure analyses add two further boundaries to any design
rule. The positive population sign survives on a volumetric basis, but the strong
coupling between density change and gravimetric response shows that pore space and
framework density are not independent scalar knobs. Across operating conditions,
pressure activation is about 4.5 times more common than pressure suppression, and
the same frozen pair can shift from neutral or unfavorable at low loading to
favorable at higher loading. Capacity gains can also oppose selectivity. These
findings align with prior evidence that methane storage is governed by coupled
porosity, density, binding, and operating-window effects
\cite{gomezgualdron2014limits,korman2020lowpressure,song2026poreengineering}, but
show that the same regime dependence persists inside chemistry- and architecture-
matched families.

We therefore propose a bounded materials-design view: pore geometry in MOFs should
be treated as a multidimensional, operating-condition-dependent regime defined by
geometry mode, geometry magnitude, coupled density context, and process objective,
rather than as a scalar property to maximize. The matched-family design strengthens
the evidential basis for that statement, but it does not convert association into a
causal deformation experiment. The chemistry representation is not atom-by-atom
identity, the largest-contrast descriptor is not an isolated causal variable, and
99.2\% of the frozen pairs are hypothetical--hypothetical. The next decisive test is
therefore targeted experimental validation in synthesized, activated, structurally
characterized matched families spanning the accessible-space and diameter regimes
identified here.

\section*{Reproducibility}

The complete analysis is organized in the public publication and reproducibility
repository at \url{https://github.com/alimardani76/isochemical_mof}. The repository
contains stage-by-stage analysis and audit scripts, frozen result tables,
figure-production code, manuscript and Supporting Information source, environment
information, and provenance manifests. Historical path-dependent scripts are
accompanied by repository reproduction instructions; post-freeze submission-stage
scripts accept explicit project/output paths.

The frozen definitions required to reproduce the reported analysis include the
recorded-chemistry key, recorded topology/dimensionality refinement, 10\% primary
geometry threshold, deterministic one-pair-per-family rule, pair orientation,
CH$_4$ endpoint at 298 K and 65 bar, median family-level estimand, and the
10,000-replicate bootstrap with seed 17. These choices should not be modified
silently when regenerating figures or downstream summaries.

Original ARC--MOF source data are not duplicated in the publication repository.
The analysis used ARC--MOF release v7, available from Zenodo at
\url{https://doi.org/10.5281/zenodo.16802743} \cite{burner2025arcmofdata},
including the charged CIF archive and the topology, dimensionality, geometry,
RAC/cluster, adsorption, and process tables described in the Supporting
Information. The complete source-to-output/code map is reported in \SItab{25},
and the final audit trail is summarized in \SIsec{21}. Required source filenames,
local hashes, repository hashes, and large-file inventories are recorded under
\texttt{data\_manifest/} and \texttt{provenance/}.

\section*{Data availability}

This study analysed publicly available ARC--MOF data from release v7, Zenodo DOI
\url{https://doi.org/10.5281/zenodo.16802743}
\cite{burner2025arcmofdata}. The upstream charged CIF archive and associated
topology, dimensionality, geometric-property, RAC/chemistry-cluster, adsorption,
and process tables are not redistributed in this repository and remain subject
to the terms of the source release.

The frozen processed outputs required to verify the reported results, together
with the analysis and independent-audit scripts, figure-generation code,
source-data inventories and hashes, environment information, manuscript and SI
source, and provenance records, are publicly available at
\url{https://github.com/alimardani76/isochemical_mof}
\cite{alimardaniabaei2026repository}. The scientific analysis and
publication-figure checkpoint used for this manuscript pass is the immutable Git
commit \texttt{88ec751e570c0519bc8f503732ccd74fc02f5f73}, available at
\url{https://github.com/alimardani76/isochemical_mof/commit/88ec751e570c0519bc8f503732ccd74fc02f5f73}.
The Supporting Information provides the complete claim--evidence and source/code
maps (Tables S22--S25) and the audit trail. No separate DOI has been minted for
the study repository at this checkpoint; the public Git commit above uniquely
identifies the version used here.

\clearpage
\bibliographystyle{unsrtnat}
\bibliography{Project17_References}

\end{document}
